\chapter{Theories and Axiomatization}
\label{ch:fol-theories}

\begin{definition}[Theory, model, satisfiability, and structure theory]
\label{def:fol-theory-satisfiability}
A \emph{theory} in \(\Lang\) is a set
\(T\subseteq\Sent(\Lang)\).  A \(\Lang\)-structure \(\Meta\) is a
\emph{model} of \(T\), written \(\Meta\entails T\), iff
\(\Meta\entails\sigma\) for every \(\sigma\in T\).  The class of models of
\(T\) is \(\Mod(T)\).

For any \(\Gamma\subseteq\Fm(\Lang)\), say that \(\Gamma\) is
\emph{satisfiable} iff there are a \(\Lang\)-structure \(\Meta\) and an
assignment \(s\) such that \(\Meta,s\entails\gamma\) for every
\(\gamma\in\Gamma\).  Because sentences are assignment-independent, a theory
is satisfiable iff it has a model.

For a \(\Lang\)-structure \(\Meta\), define its \emph{structure theory} by
\[
 \Th_{\Lang}(\Meta)
 \defeq\set{\sigma\in\Sent(\Lang):\Meta\entails\sigma}.
\]
When the signature is fixed by the surrounding statement, \(\Th(\Meta)\) is
the metalinguistic abbreviation for \(\Th_{\Lang}(\Meta)\).  A theory is an
arbitrary set of sentences, whereas a structure theory contains exactly all
sentences true in one specified structure.  Both notions are required below:
the first axiomatizes classes, and the second states elementary equivalence and
forms elementary diagrams.
\end{definition}

A theory is not required to be closed under semantic consequence.  Its
semantic closure is
\[
 \Cn(T)\defeq\set{\sigma\in\Sent(\Lang):T\entails\sigma}.
\]

\begin{definition}[Axiomatization]
A theory \(T\) \emph{axiomatizes} a class \(\mathcal K\) of
\(\Lang\)-structures iff \(\Mod(T)=\mathcal K\).  A class is
\emph{finitely axiomatizable} iff some finite theory axiomatizes it.
\end{definition}

\begin{example}[Groups]
Let \(\Lang_{\mathrm{grp}}\) contain a binary function symbol \(\cdot\), a
unary function symbol \(i\), and a constant symbol \(e\).  The following
finite theory axiomatizes groups:
\[
\begin{aligned}
&\forall x\forall y\forall z\,
  ((x\cdot y)\cdot z=x\cdot(y\cdot z)),\\
&\forall x\,(e\cdot x=x\land x\cdot e=x),\\
&\forall x\,(i(x)\cdot x=e\land x\cdot i(x)=e).
\end{aligned}
\]
The symbols are interpreted by the structure; the equations constrain those
interpretations.  Equality itself remains actual identity.
\end{example}

\begin{example}[Dense linear orders without endpoints]
In the signature containing one binary relation symbol \(<\), the following
sentences axiomatize dense linear orders without endpoints:
\[
\begin{aligned}
&\forall x\,\neg(x<x),\\
&\forall x\forall y\forall z\,
  ((x<y\land y<z)\to x<z),\\
&\forall x\forall y\,(x<y\lor x=y\lor y<x),\\
&\forall x\forall y\,(x<y\to\exists z(x<z\land z<y)),\\
&\forall x\,\exists y\exists z\,(y<x\land x<z).
\end{aligned}
\]
Each sentence contributes a distinct condition: irreflexivity, transitivity,
linearity, density, and absence of endpoints.
\end{example}

\begin{example}[Infinite structures]
For \(n\geq2\), let
\[
 \iota_n\defeq
 \exists x_0\cdots\exists x_{n-1}
 \bigwedge_{0\leq i<j<n}x_i\ne x_j,
\]
where the pairs \((i,j)\) are ordered lexicographically and the finite
conjunction uses the global right-association convention.  The notation
\(r\ne t\) abbreviates \(\neg(r=t)\).  A structure satisfies every
member of \(T_\infty=\set{\iota_n:n\geq2}\) iff its domain is
infinite.  This is an axiomatization by a theory, not by one sentence.
\end{example}

\begin{example}[Definitions inside a theory]
Suppose \(T\) is a theory in \(\Lang\), \(R\) is an \(n\)-ary relation
symbol not in \(\Lang\), and \(\varphi(x_1,\ldots,x_n)\) is an
\(\Lang\)-formula whose free variables are among the \(n\) distinct
displayed variables.  Add the \emph{defining axiom}
\[
 \delta_R\defeq
 \forall x_1\cdots\forall x_n
 \bigl(R(x_1,\ldots,x_n)\leftrightarrow
       \varphi(x_1,\ldots,x_n)\bigr).
\]
The passage from \(T\) to \(T\cup\set{\delta_R}\) is a
\emph{definitional extension}, not a new axiom about the old signature:
by \cref{thm:definitional-conservativity} below, every \(\Lang\)-model of
\(T\) has exactly one expansion satisfying \(\delta_R\), and the extended
theory has exactly the old \(\Lang\)-consequences.
\end{example}

\begin{theorem}[Conservativity of definitional extensions]
\label{thm:definitional-conservativity}
Let \(T\), \(R\), \(\varphi\), and \(\delta_R\) be as in the preceding
example, and let \(\Lang_R\) denote the signature extending \(\Lang\) by
the single relation symbol \(R\).
\begin{enumerate}
  \item Every \(\Lang\)-model \(\Meta\) of \(T\) has exactly one
        \(\Lang_R\)-expansion satisfying \(\delta_R\), and that expansion
        is a model of \(T\cup\set{\delta_R}\).
  \item The extension is \emph{conservative}: for every
        \(\sigma\in\Sent(\Lang)\),
        \(T\cup\set{\delta_R}\entails\sigma\) iff
        \(T\entails\sigma\).
\end{enumerate}
\textit{Status: complete proof.}
\end{theorem}
\begin{proof}
Both clauses use invariance under reduct
(\cref{lem:reduct-invariance}); that lemma is stated in
\cref{ch:fol-deduction}, and its proof uses only the satisfaction clauses,
not any material of this chapter.  Every member of
\(T\cup\set{\delta_R,\sigma}\) is a sentence, so satisfaction is
assignment-independent and it suffices to reason with models.

For clause 1, let \(\Meta\) be an \(\Lang\)-model of \(T\).  Since
\(\FV(\varphi)\subseteq\set{x_1,\ldots,x_n}\), formula coincidence
(\cref{lem:fol-coincidence}) shows that whether
\(\Meta,s[x_1\mapsto a_1]\cdots[x_n\mapsto a_n]\entails\varphi\) holds
depends only on the tuple \((a_1,\ldots,a_n)\), not on \(s\); define
\(R^*\subseteq(\dom\Meta)^n\) as the set of tuples for which it holds.
Now let \(\Meta'\) be any \(\Lang_R\)-expansion of \(\Meta\), and let
\(s\) be any assignment.  Unwinding the satisfaction clause for
\(\forall\) at each of the \(n\) distinct quantified variables and then
the clauses for \(\leftrightarrow\), \(\Meta'\) satisfies \(\delta_R\)
iff for every tuple \((a_1,\ldots,a_n)\in(\dom\Meta)^n\), writing
\(s'=s[x_1\mapsto a_1]\cdots[x_n\mapsto a_n]\),
\[
 \Meta',s'\entails R(x_1,\ldots,x_n)
 \quad\text{iff}\quad
 \Meta',s'\entails\varphi.
\]
The atomic clause evaluates the left side as
\((a_1,\ldots,a_n)\in R^{\Meta'}\), and because \(\varphi\) is an
\(\Lang\)-formula, invariance under reduct evaluates the right side as
\((a_1,\ldots,a_n)\in R^*\).  Hence \(\Meta'\) satisfies \(\delta_R\)
iff \(R^{\Meta'}=R^*\): the expansion interpreting \(R\) as \(R^*\) is
the unique one satisfying \(\delta_R\).  It also satisfies every member
of \(T\), because each is an \(\Lang\)-sentence true in its reduct
\(\Meta\).

For clause 2, first suppose \(T\entails\sigma\), and let \(\mathcal N\)
be any model of \(T\cup\set{\delta_R}\) in the signature \(\Lang_R\).
By invariance under reduct, \(\mathcal N\restr\Lang\) is an
\(\Lang\)-model of \(T\); hence it satisfies \(\sigma\), and invariance
under reduct returns \(\mathcal N\entails\sigma\).  Thus
\(T\cup\set{\delta_R}\entails\sigma\).  Conversely, suppose
\(T\cup\set{\delta_R}\entails\sigma\), and let \(\Meta\) be any
\(\Lang\)-model of \(T\).  Clause 1 expands \(\Meta\) to a model of
\(T\cup\set{\delta_R}\), which therefore satisfies \(\sigma\); since
\(\sigma\) is an \(\Lang\)-sentence, invariance under reduct gives
\(\Meta\entails\sigma\).  Thus \(T\entails\sigma\).
\end{proof}

Once soundness (\cref{thm:fol-soundness}) and completeness
(\cref{thm:fol-completeness}, under its Axiom-of-Choice hypothesis) are
available, clause 2 transfers to derivability: for every
\(\sigma\in\Sent(\Lang)\),
\(T\cup\set{\delta_R}\derives_{\ND_{\FOL}}\sigma\) iff
\(T\derives_{\ND_{\FOL}}\sigma\), since the left-to-right direction
composes soundness, clause 2, and completeness, and the converse holds
because a derivation from a set is a derivation from any superset.

